\section{\large {Introduction}}

This first assignment of the course essentially consists in steering a cursor with four different muscles. The group was requested to choose a person who would have worn the MCI to record a reach session. This was composed of a setup and calibration phase and a play part. In conclusion, once we had our data, the main objective was to rebuild the cursor from the raw EMG data and compare them.

More in detail, the assignment can be divided in 3 parts: 
\begin{itemize}
    \item Part 1: Band-passing the raw EMG signal in a range of 30-450Hz, then rectifying it and applying a low-pass filter at 5Hz. This is done to obtain envelopes.
    \item Part 2: Using the gains and the envelope to create the map that goes from muscle intensity to cursor velocity.
    \item Part 3: Integrating the cursor velocity into a position and animating it with the targets. It is requested to compare it with the logged cursor.
    \item Bonus part: Using our cursor to drive a simulated 2-link robot arm.
\end{itemize}

\subsection{MCI}
An MCI (Myoelectric Computer Interface) system is a hardware and software architecture used for neuro-rehabilitation purposes. The system acquires electric activity from the muscles (EMG signals, which stands for electromyographic signals) via surface sensors placed on the patient's limbs.
The signals are elaborated in real time and projected on a display. The patient uses this visual biofeedback to correct their incorrect muscular activation schemes to help their brain re-label correct movements. 



